\subsection{缺陷成像对比}
\label{sec:experiment-imaging}

为评估每种方法在全波形图像中保留缺陷相关特征的能力，在第\ref{sec:experiment-setup}节所述的含缺陷铝板试件上进行了B-scan重建实验。超声波信号在覆盖试件表面（包含机加工槽和缺口）的规则网格上采集。每种去噪方法独立应用于每个扫描位置，所得信号用于重建B-scan图像。

\textbf{对比的去噪方法。} 以256次平均作为B-scan重建的地面真值参考。每种方法的去噪输出用于形成两类图像：(1) 直接兰姆波到达的峰峰值振幅，(2) 门控波形的希尔伯特包络，采样自根据试件几何预测的理论到达时刻。提出的PMDF-Net使用训练后的模型权重通过单次前向传递应用。维纳滤波从201样本窗口平均的信号自协方差估计噪声功率谱密度。BM3D采用3.5的块匹配相似性阈值和适应1D超声波波形的硬阈值策略。DWT采用db4小波软阈值，在带噪训练集上通过五折交叉验证选择分解层数。所有参数在成像实验期间保持固定。

\textbf{评估协议。} 图像质量采用三个互补指标相对于高度平均参考B-scan进行评估。全参考指标包括信噪比改善（$\Delta$SNR，定义见第\ref{sec:metrics}节）和归一化MSE(NMSE)，在完整图像范围内逐像素计算。特征保留指标评估四个手动识别的缺陷边界附近位置的直接兰姆波到达时刻估计和峰振幅；误差报告为跨所有缺陷部位的均值$\pm$标准差。缺陷成像质量采用每个缺陷区域相对于仅噪声背景区域的对比度信噪比(CNR)来量化：
\begin{equation}
\text{CNR} = \frac{|\mu_{\text{defect}} - \mu_{\text{bkg}}|}{\sqrt{\sigma_{\text{defect}}^2 + \sigma_{\text{bkg}}^2}},
\end{equation}
其中$\mu$和$\sigma$分别表示缺陷和背景区域像素强度的均值和标准差。CNR针对每个机加工槽和缺口分别计算，然后取平均。使用Holm--Bonferroni校正的Wilcoxon符号秩检验评估跨方法指标差异的统计显著性（$p < 0.05$）。